% =========================================================
% Interval Arithmetic
% =========================================================

\subsection{Interval Arithmetic}

\begin{definitionbox}{Definition (Interval Addition)}
\begin{definition}[Interval Addition]\label{def:interval-add}
For closed intervals $[a,b]$ and $[c,d]$, their \emph{sum} is
\[
  [a,b] + [c,d] \;=\; [a+c,\; b+d].
\]
\end{definition}
\end{definitionbox}

\begin{remark*}[Standard quantified statement]
This definition fixes the named interval-arithmetic operation or relation by
the formula stated in the definition.
\end{remark*}

\begin{remark*}[Interpretation]
The formula records how interval data are transformed while preserving the
corresponding enclosure or containment meaning.
\end{remark*}

\begin{dependencies}
\begin{itemize}
  \item \hyperref[def:closed-interval]{Closed Interval}
  \item \hyperref[def:addition-on-r]{Addition on $\mathbb{R}$}
\end{itemize}
\end{dependencies}

\begin{definitionbox}{Definition (Interval Subtraction)}
\begin{definition}[Interval Subtraction]\label{def:interval-sub}
For closed intervals $[a,b]$ and $[c,d]$, their \emph{difference} is
\[
  [a,b] - [c,d] \;=\; [a-d,\; b-c].
\]
\end{definition}
\end{definitionbox}

\begin{remark*}[Standard quantified statement]
This definition fixes the named interval-arithmetic operation or relation by
the formula stated in the definition.
\end{remark*}

\begin{remark*}[Interpretation]
The formula records how interval data are transformed while preserving the
corresponding enclosure or containment meaning.
\end{remark*}

\begin{dependencies}
\begin{itemize}
  \item \hyperref[def:closed-interval]{Closed Interval}
  \item \hyperref[def:subtraction-on-r]{Subtraction on $\mathbb{R}$}
\end{itemize}
\end{dependencies}

\begin{definitionbox}{Definition (Interval Multiplication)}
\begin{definition}[Interval Multiplication]\label{def:interval-mult}
For closed intervals $[a,b]$ and $[c,d]$, their \emph{product} is
\[
  [a,b] \cdot [c,d] \;=\;
  \bigl[\min(ac, ad, bc, bd),\; \max(ac, ad, bc, bd)\bigr].
\]
\end{definition}
\end{definitionbox}

\begin{remark*}[Standard quantified statement]
This definition fixes the named interval-arithmetic operation or relation by
the formula stated in the definition.
\end{remark*}

\begin{remark*}[Interpretation]
The formula records how interval data are transformed while preserving the
corresponding enclosure or containment meaning.
\end{remark*}

\begin{dependencies}
\begin{itemize}
  \item \hyperref[def:closed-interval]{Closed Interval}
  \item \hyperref[def:multiplication-on-r]{Multiplication on $\mathbb{R}$}
  \item \hyperref[def:order-on-r]{Order on $\mathbb{R}$}
\end{itemize}
\end{dependencies}

\begin{definitionbox}{Definition (Interval Containment)}
\begin{definition}[Interval Containment]\label{def:interval-contain}
$[a,b] \subseteq [c,d]$ if and only if $c \leq a$ and $b \leq d$.
\end{definition}
\end{definitionbox}

\begin{remark*}[Standard quantified statement]
This definition fixes the named interval-arithmetic operation or relation by
the formula stated in the definition.
\end{remark*}

\begin{remark*}[Interpretation]
The formula records how interval data are transformed while preserving the
corresponding enclosure or containment meaning.
\end{remark*}

\begin{remark*}[Exposition]
Interval arithmetic provides a framework for rigorous numerical computation
with guaranteed error bounds. A real number $x \in [a,b]$ means the computed
value may lie anywhere in the interval; operations propagate these bounds.
\end{remark*}

\begin{remark*}[Exposition]
\textit{This section is a stub. A full treatment -- the sub-distributivity law,
dependency problem, applications to root-finding, and connections to
compactness -- will be developed in a future revision.}
\end{remark*}

\begin{dependencies}
\begin{itemize}
  \item \hyperref[def:closed-interval]{Closed Interval}
  \item \hyperref[def:subset]{Subset}
  \item \hyperref[def:order-on-r]{Order on $\mathbb{R}$}
\end{itemize}
\end{dependencies}
